\subsection{Stability results}

We recall the stability lemma which follow immediately from étale sheafification being a lex modality.

\begin{proposition}
TODO
\end{proposition}

\begin{lemma}\label{etale-proposition-union}
Assume given $(P_n)_{n:\N}$ a family of propositions that are étale sheaves such that $\prod_{n:\N}\, P_n\to P_{n+1}$, then $\exists(n:\N).\, P_n$ is an étale sheaf.
\end{lemma}

\begin{proof}
TODO
\end{proof}

\begin{lemma}\label{etale-set-union}
Assume given $(X_n)_{n:\N}$ a family of sets that are étale sheaves such that $\forall_{n:\N}\, X_n\subset X_{n+1}$. Then $\mathrm{colim}_{n:\N}X_n$ is an étale sheaf.
\end{lemma}

\begin{proof}
TODO
\end{proof}

\rednote{Are étale sheaf stable under sequential colimits?}


\subsection{Examples of sets that are \'etale sheaves}

\begin{proposition}
The following sets are all \'etale sheaves: 
\begin{enumerate}[(i)]
\item $\N$
\item Any scheme is an \'etale sheaf.
\item Any finitely presented modules is an \'etale sheaf.
\item Any finitely presented algebra is an \'etale sheaf.
\item TODO
\end{enumerate}
\end{proposition}

\begin{proof}
As follows:
\begin{itemize}
\item[(i)] For $\N$ it follows from \Cref{etale-set-union}. 
\item[(ii,iii,iv)] See Severi-Brauer draft. Affine schemes, f.p. algebra and f.p. and co-f.p. module immediately follow from $R$ being and \'etale sheaf and duality. For scheme we crucially use \Cref{etale-characterisation-monic}. 
\item TODO
\end{itemize}
\end{proof}


\subsection{Various descent results}

Note that we given $M$ a module that is an \'etale sheaf and $A$ an f.p.\ algebra, we don't know whether $A\otimes M$ is an \'etale sheaf. This makes the proof of descent for quasi-coherent modules much more complicated. Nevertheless we carry it, relying crucially on the characterisation of \'etale sheaf using monic unramifiable polynomials.  

\begin{lemma}\label{quasicoherent-power}
If $M$ is a quasi-coherent module and $A$ is a finitely presented algebra, then $A\otimes M$ is a quasi-coherent module.
\end{lemma}

\begin{proof}
Given $B$ a finitely presented algebra, then:
\[B\otimes A\otimes M = M^{\Spec(A)\times\Spec(B)} = (A\otimes M)^{\Spec(B)}\]
\end{proof}

\begin{lemma}\label{sheaf-unpower}
Assume $X$ a set such that $X^n$ is an \'etale sheaf for some $n\geq 1$. Then $X$ is an \'etale sheaf.
\end{lemma}

\begin{proof}
Identities in $X$ are \'etale sheaf because given $x,y:X$, we have that $(x=y) = (x=y)^n$ which is an identity type in the \'etale sheaf $X^n$. Consider the diagonal map $X\to X^n$. Its fibers are products of identity types and therefore an \'etale sheaf, its target is an \'etale sheaf therefore its source is an \'etale sheaf as needed.
\end{proof}

\begin{lemma}\label{quasi-descent-for-quasi-coherent}
Let $M$ be a module that is an \'etale sheaf and quasi-coherent up to \'etale sheafification. Then for any f.p. algebra $B$ that is finite free as a module we have that the map $B\otimes M\to M^{\Spec(B)}$ is an equivalence.
\end{lemma}

\begin{proof}
In this situation $B\otimes M$ is an \'etale sheaf as it is equal to $M^n$ as a type, so we conclude immediately from the hypothesis of quasi-coherence up to \'etale sheafification.
\end{proof}

\begin{lemma}\label{descent-quasicoherent}
Let $M$ be a module that is an \'etale sheaf. The proposition ``$M$ is quasi-coherent'' is an \'etale sheaf.
\end{lemma}

\begin{proof}
By \Cref{etale-characterisation-monic} it is enough to prove $M$ quasi-coherent assuming $P$ monic unramifiable such that writing $B=R[X]/P$ we have that $\Spec(B)\to (M\ \mathrm{q.c.})$. Under this hypothesis, assume $A$ a f.p.\ algebra, we know that $S$ implies that the map:
\[A\otimes M\to M^{\Spec(A)}\] 
is an equivalence, so that the map:
\[(A\otimes M)^{\Spec(B)}\to M^{\Spec(A)\times\Spec(B)}\]
is an equivalence. Then $(A\otimes M)^{\Spec(B)}$ is an \'etale sheaf, so that by \Cref{sheaf-unpower} we have that $A\otimes M$ is an \'etale sheaf. By \Cref{quasicoherent-power} we have that this \'etale sheaf is quasi-coherent up to \'etale sheafifcation. 

Then by \Cref{quasi-descent-for-quasi-coherent} we have that $B\otimes A\otimes M = (A\otimes M)^{\Spec(B)}$. Similarly we have that $B\otimes (M^{\Spec(A)}) = M^{\Spec(A)\times\Spec(B)}$, which all together give us an equivalence:
\[B\otimes A\otimes M \to B\otimes (M^{\Spec(A)}\]
But since $B$ is faithfully flat we conclude that the map 
\[A\otimes M\to M^{\Spec(A)}\]
is an equivalence as desired.
\end{proof}

Now that we have descent for quasi-coherent modules, we can derive further coherence much nicely, more or less directly from Lombardi-Quitt\'e. We restrict our attention to quasi-coherent modules. \rednote{it is unclear to me whether this is the right notion.}

\begin{definition}
A class of module (resp.\ algbra) $P$ is said to have algebraic \'etale descent if for all module (resp.\ algebra) $M$ that is an \'etale sheaf and quasi-coherent, we have that:
\begin{enumerate}[(a)]
\item Given $S$ in the \'etale topology such that $S\to (M\in P)$, we have that $M^S$ is in $P$.
\item Given $S$ in the \'etale topology such that $R^S\otimes M\in P$, then $M\in P$.
\end{enumerate}
\end{definition}

\begin{lemma}\label{algebraic-descent-gives-descent}
Assume given a module (resp.\ an algebra) $M$ that is an \'etale sheaf and a propriety $P$ of which has algebraic \'etale descent. Then the proposition $``M\ \mathrm{is\ quasi-coherent\ and\ has}\ P''$ is an \'etale sheaf.
\end{lemma}

\begin{proof}
We can assume $M$ quasi-coherent. Then there exists $S$ in the \'etale topology such that $S$ implies $M\in P$. By (a) we have that we have that $M^S\in P$, therefore by quasi-coherence $M\otimes R^S \in P$ so that by (c) we have that $M\in P$. 
\end{proof}

\begin{lemma}\label{lombardi-algebraic-descent}
The following classes have algebraic fppf descent.
\begin{enumerate}[(i)]
\item Finitely presented algebras.
\item Flat modules. \rednote{maybe needs f.p.}
\item Faithfully flat modules. \rednote{maybe needs f.p.}
\end{enumerate}
\end{lemma}

\begin{proof}
Condition (a) is known for finitely presented algebras. Now assume $M$ a module quasi-coherent module that is and \'etale sheaf such that there exists $S$ in the \'etale topology such that $S$ implies $M$ flat. TODO \rednote{I am having doubts about this approach}.

Condition (b) is Lombardi-Quitt\'e Theorem VIII.6.7 and VIII.6.8.
\end{proof}



\subsection{The descent compendium}

\begin{proposition}\label{proposition-all-descent}
Assume given a proposition $P$ that is an \'etale sheaf, the following propositions are \'etale sheaves:
\begin{enumerate}[(i)]
\item P is open.
\item TODO
\end{enumerate}
\end{proposition}

\begin{proof}
As follows:
\begin{itemize}
\item[(i)] See Severi-Brauer draft.
\item TODO
\end{itemize}
\end{proof}

\begin{proposition}\label{module-all-descent}
Assume given an $R$-module $M$ that is an \'etale sheaf, the following propositions are \'etale sheaves:
\begin{enumerate}[(i)]
\item $M$ is finitely presented.
\item $M$ is finitely copresented.
\item $M$ is finite free.
\item $M$ is quasi-coherent.
\item $M$ is quasi-coherent and flat. \rednote{maybe quasi-coherent needs to be replaced by f.p.?}
\item $M$ is quasi-coherent and faithfully flat. \rednote{maybe quasi-coherent needs to be replaced by f.p.?}
\item TODO
\end{enumerate}
\end{proposition}

\begin{proof}
As follows:
\begin{itemize}
\item[(i,ii,iii)] See Severi Brauer draft.
\item[(iv)] See \Cref{descent-quasicoherent}.
\item[(v,vi)] By \Cref{algebraic-descent-gives-descent} and \Cref{lombardi-algebraic-descent}.
\item TODO
\end{itemize}
\end{proof}

\begin{proposition}\label{algebra-all-descent}
Assume given an algebra $A$ that is an \'etale sheaf, the following propositions are \'etale sheaves:
\begin{enumerate}[(i)]
\item $A$ is finitely presented.
\item TODO
\end{enumerate}
\end{proposition}

\begin{proof}
As follows:
\begin{itemize}
\item[(i)] By \Cref{algebraic-descent-gives-descent} and \Cref{lombardi-algebraic-descent}.
\item TODO
\end{itemize}
\end{proof}

\begin{proposition}\label{types-all-descent}
Assume given an \'etale sheaf $X$, the following propositions are \'etale sheaves:
\begin{enumerate}[(i)]
\item $X$ is an affine scheme.
\item $X$ is formally \'etale.
\item $X$ is formally unramified.
\item $X$ is a Deligne-Mumford stack.
\item TODO
\end{enumerate}
\end{proposition}

\begin{proof}
As follows:
\begin{itemize}
\item[(i)] Dual of \Cref{algebra-descent}
\item[(ii,iii)] Just from writing the definitions.
\item[(iv)] From Tim's master thesis.
\item TODO
\end{itemize}
\end{proof}

\begin{remark}
Beware that \'etale descent for schemes fails.
\end{remark}





